\section{Discussion}

\textsc{StructIRES} separates two related but distinct questions: whether a
sequence is IRES-like in a defined identification benchmark, and whether a
candidate can be optimized without losing the predicted structural state of an
experimentally supported parent. The completed \textsc{StructIRES-Rank}
component turns an IRES-recognition score into one component of a controlled
multi-objective selection procedure. Under a fixed proposal budget, it selects
candidates that retain parent energy and ensemble state more faithfully while
improving a held-out direct-RNA computational proxy. The IAPV case further
shows how direct-RNA mutational knowledge can be added transparently as a local
edit-risk guardrail.

The structure-aware recognition component is intentionally not claimed as a
completed performance gain in this draft. It becomes a contribution only if a
validation-locked run improves held-out ranking beyond the sequence baseline;
otherwise the paper remains a constrained-design study with transparent
sequence baselines.

The study has limitations. Secondary-structure predictions do not represent
all cellular or circular RNA effects, and the direct-RNA proxy is not a
functional assay. We therefore make no claim of experimentally verified IRES
activity for selected variants. Future work should evaluate the frozen
shortlists prospectively in matched direct-RNA settings and extend the
assay-aware framework to cargo and cell-context robustness.
